\section{先写下一页}
\subsection{把作品带进来}
项目首页左侧是书架，选中一本，右侧可以翻阅最近编译的 PDF；点击\menu{进入项目}开始编辑。书架可以独立滚动，项目旁的菜单用于改名或删除。

\menu{上传／导入}可以带入单个文件、多个文件、文件夹或 ZIP 压缩包。导入文件夹时保留内部目录，ZIP 会展开为项目文件。已有项目里也可以继续上传。文件夹选择需要浏览器支持，空目录不会随文件夹导入；ZIP 请单独选择。遇到同名路径时先改名再导入，已有文件不会被静默覆盖。

左侧文件列表支持新建文件、目录、改名和拖动移动。开启\menu{批量选择}后可以导出、移动或删除多个文件；处理结果会逐项显示。删除图片等附件前，先导出一份副本，文本历史不能代替二进制附件备份。

\subsection{源文件与 PDF}
编辑区里的 \texttt{main.tex} 是这份手册的主文件，正文通过 \verb|\input| 引入 \texttt{chapters/} 下的各章。点击\menu{编译}后，右侧 PDF 展示整个项目的结果；单独打开某一章，不会把编译入口自动换成那一章。

需要换入口时，在项目设置中选择\menu{主文件}。同一处可以选择编译引擎：中文或自定义字体项目通常用 XeLaTeX；已有模板则先沿用作者的设置。这个手册已经选择了 XeLaTeX，直接编译即可。

\tryit{打开本章文件，把下面一句改成你想写的题目，等自动保存完成后点击编译。}
\begin{quote}我正在写一份关于人与城市的文档。\end{quote}

\subsection{把阅读环境调舒服}
\menu{设置 → 外观}里可以选一整套日夜配色，切换浅色、深色或跟随系统，也可以调整界面字体、书名字体和 Prims 字标。它们不会改变论文的 LaTeX 字体或排版。

PDF 区域的外观下拉可以选择阅读效果，包括暖黄或浅绿护眼色。阅读滤镜只影响屏幕显示，下载出来的 PDF 仍保留文档原本的颜色。
